\documentclass[10pt,a4paper]{article}
\usepackage[utf8]{inputenc}
\usepackage[T1]{fontenc}
\usepackage[margin=1.8cm,top=2.6cm,bottom=2.2cm]{geometry}
\usepackage{parskip}
\usepackage{microtype}
\usepackage{amsmath,amssymb,mathtools}
\usepackage{array,booktabs,tabularx,multirow}
\usepackage[table]{xcolor}
\usepackage{graphicx,float}
\usepackage{titlesec}
\usepackage{enumitem}
\usepackage{fancyhdr}

\renewcommand\thesection{\Roman{section}}
\numberwithin{equation}{section}

\titleformat{\section}{\large\bfseries}{\thesection.\;}{0em}{}[\titlerule]
\titleformat{\subsection}{\normalsize\bfseries}{}{0em}{}
\titleformat{\subsubsection}{\normalsize\itshape}{}{0em}{}
\titlespacing{\section}{0pt}{1.6ex plus .2ex}{0.6ex}
\titlespacing{\subsection}{0pt}{1.0ex}{0.3ex}

\pagestyle{fancy}
\fancyhf{}
\setlength{\headheight}{15pt}
\fancyhead[L]{\small\textit{Layer~12: SWV Peak Extraction}}
\fancyhead[R]{\small Aaran Nihalani $\bullet$ Eton iGEM Dry Lab}
\fancyfoot[C]{\small\thepage}
\renewcommand{\headrulewidth}{0.35pt}

\newcommand{\cit}[1]{[#1]}
\newcommand{\kB}{k_{\text{B}}}
\newcommand{\kET}{k_{\text{ET}}^{0}}
\newcommand{\Eo}{E^{\circ\prime}}
\newcommand{\Esw}{E_{\text{sw}}}
\newcommand{\Apk}{A_{\text{peak}}}
\newcommand{\Aarea}{A_{\text{area}}}
\newcommand{\Epk}{E_{\text{peak}}}
\newcommand{\DI}{\Delta I}

\begin{document}

%=======================================================================
\begin{center}
{\LARGE\bfseries Layer~12: SWV Peak Extraction and Signal Processing}\\[0.4em]
{\large From Raw Voltammogram to Peak Areas with Quantified Uncertainty}\\[0.25em]
{\small\color{gray} Eton iGEM 2026, Dry Lab \quad$|$\quad
Layer~12 of 17-Layer Biosensor Pipeline \quad$|$\quad
\texttt{peak\_extraction.py}}
\end{center}
\vspace{0.4em}\hrule\vspace{0.7em}

\noindent
Layer~11 produces a theoretical net SWV current trace $\DI(E)$.
The instrument measures a noisy, baseline-corrupted version of this trace.
Layer~12 takes the raw measured array $\{(E_n, \DI_n)\}$ and extracts
two numbers: the MB peak area $A_{\text{area}}^{\text{MB}}$ and the Fc peak area
$A_{\text{area}}^{\text{Fc}}$, each with a propagated uncertainty $\sigma$.
These two quantities are the sole inputs to Layer~13.

The importance of this layer cannot be overstated. No matter how accurate
the physical models in Layers~0--11 are, a systematic error introduced here
(wrong baseline, clipped peak, noise spike misidentified as signal) corrupts
every downstream calculation: the ratiometric signal $F_k$, the apparent
rate $\kappa_k$, and the final enzyme concentration. The pipeline therefore
applies four sequential operations: noise filtering (Section~\ref{sec:sg}),
background subtraction (Section~\ref{sec:baseline}), peak fitting
(Section~\ref{sec:fit}), and quality diagnostics (Section~\ref{sec:qc}).

%=======================================================================
\section{Physical Shape of the SWV Peak}
\label{sec:shape}

Before fitting a model to the data, it is essential to derive the expected
peak shape from first principles so that the choice of fitting function can
be justified rather than assumed.

\subsection{Fermi-Dirac occupancy and the sech$^2$ peak}

From Layer~11, the net SWV current at potential step $n$ is
$\DI_n = I_\text{fwd}(E_n + \Esw) - I_\text{rev}(E_n - \Esw)$. In the
reversible limit ($\kappa_\text{SW} \gg 1$), each half-cycle reaches Nernst
equilibrium before current is sampled. The equilibrium surface coverage of
oxidised MB at potential $E$ is set by the Fermi-Dirac occupancy function:
\begin{equation}
  \Gamma_\text{ox}^\infty(E) = \frac{\Gamma_\text{intact}}{1 +
  \exp\!\bigl(n_e F(E - \Eo_\text{MB})/RT\bigr)}
  \label{eq:fermi}
\end{equation}
This equation states that at potentials much more positive than $\Eo_\text{MB}$,
MB is predominantly oxidised ($\Gamma_\text{ox} \to \Gamma_\text{intact}$);
at potentials much more negative, it is predominantly reduced
($\Gamma_\text{ox} \to 0$). The transition is centred at $E = \Eo_\text{MB}$
and its sharpness is set by $RT/(n_e F)$: lower temperature or more electrons
transferred gives a sharper transition.

The net SWV current is proportional to the difference in equilibrium coverage
between the two half-cycle potentials:
\begin{equation}
  \DI(E) \propto \Gamma_\text{ox}^\infty(E + \Esw)
  - \Gamma_\text{ox}^\infty(E - \Esw)
  \label{eq:DI_diff}
\end{equation}
For small $\Esw$ relative to $RT/(n_e F)$, this difference approximates the
derivative of $\Gamma_\text{ox}^\infty$ with respect to $E$, multiplied by
$2\Esw$. Computing that derivative from Eq.~\eqref{eq:fermi}:
\begin{equation}
  \frac{d\Gamma_\text{ox}^\infty}{dE}
  = -\frac{n_e F \Gamma_\text{intact}}{4RT}\,
  \text{sech}^2\!\!\left(\frac{n_e F (E - \Eo_\text{MB})}{2RT}\right)
  \label{eq:deriv}
\end{equation}
The net SWV current therefore has the shape:
\begin{equation}
  \boxed{
  \DI(E) \propto \Gamma_\text{intact} \cdot
  \text{sech}^2\!\!\left(\frac{n_e F(E - \Eo_\text{MB})}{2RT}\right)
  }
  \label{eq:sech2}
\end{equation}
This result, confirmed by Mirceski et~al.~\cit{1} for surface-confined
reversible systems, tells us three things directly. First, the peak is
\emph{proportional to} $\Gamma_\text{intact}$: it is a direct measure of the
remaining intact probe density after cleavage, which is the quantity Layer~13
must extract. Second, the peak is centred exactly at $\Eo_\text{MB}$, allowing
a cross-check against the pH-corrected value from Layer~10. Third, the peak
shape is entirely determined by the thermodynamic ratio $RT/(n_e F)$: it is
not a free parameter of the fitting model.

\subsection{Full-width at half-maximum and the Gaussian approximation}

The FWHM of Eq.~\eqref{eq:sech2} is obtained by solving $\text{sech}^2(x) = 1/2$:
\begin{equation}
  \text{FWHM} = \frac{2 \times 2\ln(1+\sqrt{2}) \times RT}{n_e F}
  = \frac{2 \times 1.763 \times RT}{n_e F}
  = \frac{3.526 \times 8.314 \times 310}{2 \times 96485}
  \approx 47\;\text{mV}
  \label{eq:fwhm}
\end{equation}
at $T = 310$~K, $n_e = 2$. The corresponding Gaussian standard deviation is:
\begin{equation}
  \sigma_E^\text{pred} = \frac{\text{FWHM}}{2\sqrt{2\ln 2}}
  = \frac{\text{FWHM}}{2.355}
  \approx 20\;\text{mV}
  \label{eq:sigpred}
\end{equation}
The Gaussian function $\exp(-x^2/2)$ approximates $\text{sech}^2(ax)$ to
within 5\% for $|x| \leq 1.5\sigma$ when the Gaussian $\sigma$ is matched to
the sech$^2$ FWHM via Eq.~\eqref{eq:sigpred}. Beyond $\pm 1.5\sigma$, the
sech$^2$ has heavier tails than the Gaussian; this 5\% discrepancy in the
wings contributes approximately $0.3\%$ error to the integrated peak area.
The Gaussian fit is therefore an excellent approximation for area extraction
and is preferred over sech$^2$ for numerical stability.

In the quasi-reversible regime ($\kappa_\text{SW} \sim 1$), the peak broadens
and its height decreases compared to the reversible limit. Mirceski et~al.
\cit{1} show that the peak shape remains approximately Gaussian in the
quasi-reversible regime, with $\sigma_E$ increasing beyond the reversible
prediction. A measured $\sigma_E > 1.5 \times \sigma_E^\text{pred}$ is
therefore a sign that the system is in the quasi-reversible regime at the
chosen $f_\text{SWV}$, not of heterogeneous SAM chemistry.

%=======================================================================
\section{Savitzky-Golay Noise Filtering}
\label{sec:sg}

\subsection{The convolution equation}

The potentiostat samples the current at $N \approx 650$ discrete potential
steps (for a $-0.65$ to $+0.65$~V scan at 2~mV step, from Baker~2006~\cit{3}
and Lubin~2010~\cit{4}). Each sampled value $\DI_n$ carries per-point noise
$\delta_n$ with standard deviation $\sigma_\text{noise}$ (from the
transimpedance amplifier and analogue-to-digital converter, Layer~14).
The Savitzky-Golay (SG) filter~\cit{2} replaces each point with a
weighted sum of its neighbours:
\begin{equation}
  \hat{I}_n = \sum_{k=-m}^{m} c_k\, \DI_{n+k}
  \label{eq:sg}
\end{equation}
where $m$ is the half-window width and the weights $\{c_k\}$ are
determined by fitting a polynomial of degree $p$ to the $2m+1$
points in the window and reading off the polynomial value at the
centre point. The key property of these weights is that they satisfy:
\begin{equation}
  \sum_{k=-m}^{m} c_k = 1, \qquad
  \sum_{k=-m}^{m} c_k\, k^j = 0 \quad \text{for } j = 1, \ldots, p
  \label{eq:sgconds}
\end{equation}
The first condition in Eq.~\eqref{eq:sgconds} ensures that a constant signal
passes through the filter unchanged (no DC bias). The second condition
ensures that polynomial signals of degree up to $p$ also pass unchanged,
which is why the filter preserves curved features such as peak tops.

\subsection{Why SG outperforms a running average for peaked data}

A running average (boxcar filter) is the special case $c_k = 1/(2m+1)$
for all $k$, corresponding to $p = 0$ (degree-0 polynomial, a flat line).
A flat-line fit is a poor model for a peak: it systematically
underestimates the signal at the peak apex and overestimates it on the
flanks, broadening the peak in proportion to the window width.
The SG filter with $p = 3$ fits a cubic polynomial locally, which can
reproduce the curvature at the peak top without distorting it, as long
as the peak is wider than the window.

\subsection{Noise reduction factor}

Because the weights $\{c_k\}$ are fixed (not random), and because the noise
on each point $\delta_n$ is uncorrelated (white noise from the ADC and TIA),
the noise on the smoothed point $\hat{I}_n$ is:
\begin{equation}
  \sigma_{\hat{I}} = \sigma_\text{noise}\,\sqrt{\sum_{k} c_k^2}
  \label{eq:sgnoise}
\end{equation}
For the symmetric SG filter with $m = 5$ (11 points), $p = 3$, the tabulated
sum $\sum_k c_k^2 \approx 0.233$~\cit{2}, giving:
\begin{equation}
  \sigma_{\hat{I}} = \sigma_\text{noise} \times \sqrt{0.233} \approx 0.48\,\sigma_\text{noise}
  \label{eq:sgnoise_val}
\end{equation}
The filter reduces per-point noise by a factor of $\approx 2.1\times$
relative to the unsmoothed data. This is meaningful: the per-point SNR
on the MB peak is approximately $0.27\,\mu\text{A} / \sigma_\text{noise}$;
doubling the per-point SNR halves the uncertainty on the fitted peak
amplitude.

\subsection{Parameter choice}

The window half-width $m = 5$ (11 points, 22~mV) is chosen to be well below
the peak FWHM of 47~mV (Eq.~\ref{eq:fwhm}), ensuring no significant peak
broadening. The polynomial degree $p = 3$ is the minimum needed to capture
the curvature at the peak apex (a quadratic feature); degree 5 would give
marginal improvement in peak preservation at the cost of reduced noise
suppression. These choices are consistent with the recommendations of
Savitzky and Golay~\cit{2} for spectroscopic peaks with similar
width-to-window ratios.

%=======================================================================
\section{Baseline Subtraction}
\label{sec:baseline}

\subsection{Origin of the baseline}

Even after SWV differential cancellation (Layer~11, Section~VII), a
non-zero slowly varying baseline remains in the measured trace for three
reasons. First, the capacitive transient is not perfectly symmetric between
forward and reverse half-cycles when the potential step $\Delta E_s$ is
non-zero, leaving a residual proportional to $dC_\text{dl}/dE \times
\Delta E_s$. Second, the O$_2$ reduction wave contributes a small monotone
cathodic current across the scan (Layer~11, Section~VIII). Third, unresolved
faradaic contributions from adventitious species in the sample matrix may
add broad features.

The baseline is not a physical signal of interest; it is an additive
artefact that, if not removed, biases the peak area. A 1~nA baseline
elevation across a 200~mV peak window adds approximately
$1\;\text{nA} \times 0.2\;\text{V} = 0.2\;\text{nA}\cdot\text{V}$
to the fitted area, comparable to the area of the MB peak itself at low
probe coverage.

\subsection{Peak windows and anchor regions}

Two \textit{peak windows} are defined as the voltage ranges expected to
contain Faradaic current. Everything outside these windows is treated as
baseline and used to anchor the baseline fit.

The window boundaries are set at $\pm 3\sigma_E^\text{pred}$ from the
expected peak centres:
\begin{align}
  \text{MB window:}&\quad
  [\Eo_\text{MB}(\text{pH}) - 3\sigma_E^\text{pred},\;
   \Eo_\text{MB}(\text{pH}) + 3\sigma_E^\text{pred}]
  = [\Eo_\text{MB} - 60\;\text{mV},\;\Eo_\text{MB} + 60\;\text{mV}]
  \label{eq:mbwindow}\\[0.3em]
  \text{Fc window:}&\quad
  [\Eo_\text{Fc} - 3\sigma_E^\text{pred,Fc},\;
   \Eo_\text{Fc} + 3\sigma_E^\text{pred,Fc}]
  \approx [\Eo_\text{Fc} - 90\;\text{mV},\;\Eo_\text{Fc} + 90\;\text{mV}]
  \label{eq:fcwindow}
\end{align}
The $3\sigma$ guard band ensures that less than 0.3\% of the Gaussian peak
area lies outside the window and therefore the anchor points carry
essentially no Faradaic signal. The Fc window is wider (90~mV vs 60~mV)
because the Fc peak, with $n_e = 1$, has FWHM $\approx 94$~mV at 37$^\circ$C,
giving $\sigma_E^\text{pred,Fc} \approx 40$~mV.

\subsection{Baseline model and fit}

A cubic polynomial is fitted by least squares to the $N_\text{base}$ points
in the anchor regions. The polynomial is then evaluated across the full
scan range and subtracted from the smoothed data. The baseline-subtracted
current at step $n$ is:
\begin{equation}
  \DI_n^\text{corr} = \hat{I}_n - \bigl(b_0 + b_1 E_n + b_2 E_n^2 + b_3 E_n^3\bigr)
  \label{eq:baseline}
\end{equation}
The cubic polynomial captures the slowly varying capacitive slope and any
mild nonlinearity; a linear fit would suffice for a flat baseline but
fails if the O$_2$ wave tail introduces curvature at the cathodic end of
the scan.

\subsection{Quality check on baseline fit}

The RMS residual of the polynomial fit over the anchor points is compared
to the expected noise level $\sigma_\text{noise}$:
\begin{equation}
  Q_\text{base} = \frac{1}{N_\text{base}}\sum_{n \in \text{anchor}}
  \bigl(\hat{I}_n - \hat{b}(E_n)\bigr)^2 / \sigma_\text{noise}^2
  \label{eq:Qbase}
\end{equation}
$Q_\text{base}$ follows a $\chi^2$ distribution with $N_\text{base} - 4$
degrees of freedom under the null hypothesis that the baseline is truly
cubic. If $Q_\text{base} > 5$, the baseline is non-polynomial, indicating
possible electrode fouling, an unresolved interfering wave, or a bubble at
the electrode surface. The scan is flagged and excluded from averaging.

%=======================================================================
\section{Gaussian Peak Fitting and Area Extraction}
\label{sec:fit}

\subsection{The Gaussian model}

After baseline subtraction, the baseline-corrected current in each peak
window is fitted to a Gaussian:
\begin{equation}
  \DI^\text{corr}(E) = \Apk\,\exp\!\!\left(-\frac{(E - \Epk)^2}{2\sigma_E^2}\right)
  \label{eq:gauss}
\end{equation}
This model has three free parameters. $\Apk$ (nA) is the peak
amplitude: it is proportional to $\Gamma_\text{intact}$ (the quantity
Layer~13 needs) but also to $f(\kappa_\text{SW})$, the kinetic factor
from Layer~11 that depends on $\kET$ and the scan frequency $f_\text{SWV}$.
Because $\kET$ varies between electrodes (different probe lengths, different
tether conformations), $\Apk$ alone is not a reliable measure of
$\Gamma_\text{intact}$. $\Epk$ (V) is the measured peak centre: it should
agree with the Layer~10 prediction $\Eo_\text{MB}(\text{pH})$ within a
tolerance set by the on-chip reference electrode drift (typically
$\pm 25$~mV; Baker~2006~\cit{3}, Lubin~2010~\cit{4}). A shift outside this
tolerance indicates reference electrode failure, sample pH far from
expectation, or a misidentified peak. $\sigma_E$ (V) is the peak width:
its predicted value is $\sigma_E^\text{pred} \approx 20$~mV
(Eq.~\ref{eq:sigpred}); deviations carry diagnostic information (see
Section~\ref{sec:qc}).

\subsection{Why peak area rather than peak height}

The integrated area under the Gaussian is:
\begin{equation}
  \Aarea = \int_{-\infty}^{+\infty} \Apk\,
  \exp\!\!\left(-\frac{(E-\Epk)^2}{2\sigma_E^2}\right)\!dE
  = \Apk\,\sigma_E\,\sqrt{2\pi}
  \label{eq:area}
\end{equation}
This equation says that the area is the product of the peak height $\Apk$
and the width $\sigma_E$, scaled by $\sqrt{2\pi}$. The key physical point
is that as the system moves from the reversible to the quasi-reversible
regime (lower $f_\text{SWV}$, slower $\kET$), the peak becomes shorter and
broader simultaneously: $\Apk$ decreases while $\sigma_E$ increases, with
the product $\Apk \times \sigma_E$ approximately conserved. This charge
conservation follows from the fact that the total charge passed during a
sweep through the peak is $Q = n_e F A_\text{true} \Gamma_\text{intact}$,
which depends only on the physical state of the electrode (probe coverage
and electrode area) and not on the scan kinetics.

The amplitude $\Apk$ therefore conflates the kinetic regime with the surface
coverage, whereas the area $\Aarea$ is approximately proportional to
$\Gamma_\text{intact}$ alone. White and Plaxco~\cit{5} confirm
experimentally that the SWV signal gain (effectively the peak area) is
the more robust metric for quantifying surface coverage across a range of
$f_\text{SWV}$ values. We therefore use $\Aarea$ as the primary output
of this layer, not $\Apk$.

\subsection{Nonlinear least-squares fitting and the covariance matrix}

The three parameters $(\Apk, \Epk, \sigma_E)$ are determined by minimising
the sum of squared residuals between the data and Eq.~\eqref{eq:gauss}.
The fitting algorithm (Levenberg-Marquardt, as implemented in
\texttt{scipy.optimize.curve\_fit}) provides, in addition to the best-fit
parameters, the $3\times3$ covariance matrix $\Sigma$. The diagonal
elements give the variances: $\sigma_{\Apk}^2 = \Sigma_{11}$,
$\sigma_{\Epk}^2 = \Sigma_{22}$, $\sigma_{\sigma_E}^2 = \Sigma_{33}$.
The off-diagonal element $\Sigma_{13}$ gives the covariance between
$\Apk$ and $\sigma_E$.

This covariance is physically important and must not be ignored. When a peak
is slightly wider than the model predicts, the least-squares algorithm
responds by reducing $\Apk$ and increasing $\sigma_E$, keeping their
product approximately constant. The amplitude and width are therefore
negatively correlated: $\Sigma_{13} < 0$. Propagating uncertainty in
$\Aarea = \Apk\,\sigma_E\,\sqrt{2\pi}$ while ignoring this covariance
overestimates $\sigma_{\Aarea}$. The correct propagation is:
\begin{equation}
  \sigma_{\Aarea}^2 = 2\pi\!\left[
    \sigma_E^2\,\sigma_{\Apk}^2
    + \Apk^2\,\sigma_{\sigma_E}^2
    + 2\,\Apk\,\sigma_E\,\Sigma_{13}
  \right]
  \label{eq:sigarea}
\end{equation}
The third term (the covariance contribution) is negative, so ignoring it
gives a conservative (upper-bound) estimate of $\sigma_{\Aarea}$.

\subsection{Initial guesses and fitting bounds}

The fitting algorithm requires initial guesses for the three parameters.
These are obtained from the smoothed, baseline-subtracted data: the initial
$\Apk$ is the maximum value within the peak window, $\Epk$ is the potential
at that maximum, and $\sigma_E$ is set to the predicted value
$\sigma_E^\text{pred}$. Bounds are imposed:
$\Apk \in [0, 2\max(\DI^\text{corr})]$,
$\Epk \in [\Eo_\text{MB}(\text{pH}) \pm 75\;\text{mV}]$,
$\sigma_E \in [0.5\sigma_E^\text{pred}, 5\sigma_E^\text{pred}]$.
These bounds prevent the algorithm from converging to unphysical solutions
(negative amplitude, unreasonably sharp or broad peaks).

%=======================================================================
\section{Two-Gaussian Simultaneous Fit}
\label{sec:twogauss}

When both MB and Fc peaks are extracted from the same scan, they are fitted
simultaneously with a shared baseline model:
\begin{equation}
  \DI(E) = \Apk^\text{MB}\,G_\text{MB}(E)
  + \Apk^\text{Fc}\,G_\text{Fc}(E)
  + b_0 + b_1 E
  \label{eq:twogauss}
\end{equation}
where $G_j(E) = \exp(-(E-\Epk^j)^2/(2\sigma_E^{j\,2}))$. The simultaneous
fit prevents the following bias: if the two peaks are fitted sequentially
and the first fit has slightly wrong wings, the residual from the first peak
is absorbed into the baseline of the second, systematically biasing the
second area.

The two peaks are resolvable as independent Gaussians when their separation
exceeds $3(\sigma_E^\text{MB} + \sigma_E^\text{Fc})$. For
$\sigma_E^\text{MB} \approx 20$~mV and $\sigma_E^\text{Fc} \approx 40$~mV,
the resolvability criterion requires $|\Epk^\text{Fc} - \Epk^\text{MB}|
> 180$~mV. The expected separation is approximately 690~mV (Layer~10,
Section~VIII), comfortably exceeding this threshold. The two-Gaussian model
is therefore never degenerate for this system.

%=======================================================================
\section{Quality Diagnostics}
\label{sec:qc}

Every extracted peak is subjected to the following automated diagnostics.
Each diagnostic flag has a specific physical interpretation.

\begin{table}[H]
\centering
\small\renewcommand{\arraystretch}{1.4}
\begin{tabularx}{\linewidth}{@{}p{3.5cm}p{2.5cm}p{2.5cm}X@{}}
\toprule
Diagnostic & Flag condition & Expected value & Physical interpretation \\
\midrule
Peak width $\sigma_E^\text{MB}$ & $> 2\sigma_E^\text{pred}$ & $20$~mV & SAM is heterogeneous (two populations of probe with different $\Eo_\text{MB}$), or dissolved MB is contributing a broader diffusion-controlled peak \\[0.5em]
Peak width $\sigma_E^\text{MB}$ & $< 0.5\sigma_E^\text{pred}$ & $20$~mV & Noise spike misidentified as peak; or the peak is not a Faradaic signal \\[0.5em]
Peak position $|\Epk^\text{MB} - \Eo_\text{MB}(\text{pH})|$ & $> 25$~mV & $<5$~mV & Reference electrode drift, sample pH far from the assumed value, or a different redox-active species is being fitted \\[0.5em]
Fit $\chi^2/\nu$ & $> 5$ & $\approx 1$ & Non-Gaussian peak shape; the sech$^2$ vs Gaussian discrepancy is only 5\% at the wings, so a value $> 5$ indicates a real non-Gaussianity (asymmetric peak from two probe populations, or double-MB peak) \\[0.5em]
Fc amplitude stability $|\Apk^\text{Fc}(i) - \bar{\Apk}^\text{Fc}|/\bar{\Apk}^\text{Fc}$ & $> 10\%$ between scans & $< 2\%$ & Fc SAM is damaged (Fc molecules desorbing) or a bubble is intermittently blocking part of the electrode \\[0.5em]
Signal below LOD: $\Apk < 3\sigma_\text{noise}/\sqrt{N_\text{pts}}$ & (triggered) & Not triggered & Probe coverage is too low to detect, or signal is lost through complete cleavage; Layer~13 reports $A_\text{area} = 0 \pm A_\text{LOD}$ \\[0.5em]
\bottomrule
\end{tabularx}
\caption{Automated quality diagnostics applied to every SWV peak extraction.
$\nu$ is the number of degrees of freedom in the Gaussian fit
($N_\text{pts} - 3$, where $N_\text{pts}$ is the number of points in the
peak window); $\sigma_\text{noise}$ is the per-point noise from Layer~14.}
\label{tab:qc}
\end{table}

%=======================================================================
\section{Edge Cases}
\label{sec:edge}

\subsection{Double MB peak}

If two distinct MB populations exist on the electrode (e.g., intact probe
with MB at the full tether distance, and a cleaved stub with MB at a shorter
distance), the SWV voltammogram may show two partially overlapping MB peaks
at slightly different formal potentials. The single-Gaussian fit would
return a large $\sigma_E$ and a high $\chi^2/\nu$. The diagnostic table
(Section~\ref{sec:qc}) flags this case, and the pipeline switches to a
two-Gaussian fit within the MB window. The two fitted areas are reported
separately; Layer~13 uses only the intact-probe peak (the peak at
$\Epk \approx \Eo_\text{MB}(\text{pH})$).

\subsection{Signal clipping}

If $\Apk > 0.9 \times I_\text{range,TIA}$, the transimpedance amplifier has
saturated and the peak top is clipped. The fitted Gaussian amplitude is
unreliable. This flag triggers a recommendation to reduce the probe coverage
$\Gamma_0$ or the TIA gain.

\subsection{Signal below detection limit}

If the peak amplitude is not statistically distinguishable from the noise
(see Table~\ref{tab:qc}, last row), Layer~13 is not called for this channel.
The result is reported as a censored measurement: $A_\text{area} < A_\text{LOD}$.

%=======================================================================
\section{Frequency Diagnostic Implementation}
\label{sec:freqdiag}

From Layer~11 (Section~XI), the slope of $\log\Aarea$ versus $\log f_\text{SWV}$
distinguishes surface-confined from diffusion-controlled redox. This
diagnostic is run once per electrode batch (not on every cartridge):

\begin{equation}
  \frac{d\log\Aarea}{d\log f_\text{SWV}} =
  \begin{cases}
    1 & \text{reversible surface-confined (}f \ll k^0/2\text{)}\\
    0\text{--}1 & \text{quasi-reversible surface-confined}\\
    0.5 & \text{diffusion-controlled (dissolved MB)}
  \end{cases}
  \label{eq:freqslope}
\end{equation}
The area $\Aarea$ is extracted by the Gaussian fit at each frequency using
the same pipeline described in this layer. The slope of the log-log plot
confirms which regime the electrode is operating in and validates the
Layer~11 model.

%=======================================================================
\section{Layer Outputs}
\label{sec:outputs}

\begin{table}[H]
\centering
\small\renewcommand{\arraystretch}{1.3}
\begin{tabularx}{\linewidth}{@{}lXlll@{}}
\toprule
Output & Physical meaning & Typical value & Uncertainty & Feeds \\
\midrule
$\Aarea^\text{MB}$ & Integrated MB peak area & 0.05--5~nA\,V & $\sigma_{\Aarea}$ (Eq.~\ref{eq:sigarea}) & Layer~13 \\
$\Aarea^\text{Fc}$ & Integrated Fc peak area & 0.1--10~nA\,V & $\sigma_{\Aarea}$ & Layer~13 \\
$\Epk^\text{MB}$ & Measured MB formal potential & $-0.29 \pm 0.06$~V vs on-chip RE & $\sigma_{\Epk}$ & Layer~10 cross-check \\
$\sigma_E^\text{MB}$ & MB peak width & $\approx 20$~mV & $\sigma_{\sigma_E}$ & Quality diagnostic \\
Quality flags & Boolean flags from Table~\ref{tab:qc} & All clear & --- & Data pipeline \\
\bottomrule
\end{tabularx}
\caption{Outputs of Layer~12 passed to downstream layers. All uncertainties
are at the $1\sigma$ level and come from the nonlinear least-squares
covariance matrix (Eq.~\ref{eq:sigarea}).}
\label{tab:outputs}
\end{table}

%=======================================================================
\section{Limitations}
\label{sec:limits}

\begin{enumerate}[leftmargin=*]

\item \textbf{The Gaussian is an approximation to sech$^2$.} The systematic
  error in the fitted area from the sech$^2$/Gaussian discrepancy is
  $\approx 0.3\%$ (Section~\ref{sec:shape}). This is negligible compared
  to other uncertainties, but the numerical implementation in
  \texttt{peak\_extraction.py} can optionally use the exact sech$^2$ model
  for verification.

\item \textbf{The single-Gaussian fit assumes one probe population.} If the
  SAM contains probes with a broad distribution of $\Eo_\text{MB}$ values
  (e.g., from heterogeneous PEG linker lengths), the true peak is a
  superposition of Gaussians at different centres. The fitted single Gaussian
  overestimates $\sigma_E$ and produces a biased area. The frequency
  diagnostic (Section~\ref{sec:freqdiag}) can distinguish this from
  quasi-reversible broadening: kinetic broadening increases with $f$,
  whereas heterogeneous broadening is frequency-independent.

\item \textbf{The cubic baseline polynomial fails near the O$_2$ reduction
  wave.} If the cathodic scan limit is close to the O$_2$ wave (at
  approximately $-0.5$ to $-0.6$~V vs Ag/AgCl in aerated buffer), the wave
  tail introduces a nonlinear feature at the MB baseline anchor region. In
  this case, a spline through more anchor points, or a separate exponential
  term for the O$_2$ contribution, must be added to Eq.~\eqref{eq:baseline}.

\item \textbf{The Fc peak width prediction assumes $n_e = 1$ and the
  reversible limit.} If Fc is at a distance where $k^{0}_{\text{ET,Fc}}$ is not
  much larger than $2f_\text{SWV}$, the Fc peak will broaden and its area
  will drop. Section~\ref{sec:fit} applies the same area-over-height
  argument to Fc as to MB, so the area is still the preferred metric; but
  if the Fc peak broadens enough to overlap with the MB peak tail, the
  two-Gaussian simultaneous fit (Section~\ref{sec:twogauss}) is required.

\end{enumerate}

%=======================================================================
\section*{References}

\noindent All papers in \texttt{electrochemical\_device/papers/layer\_12/}.

\begin{enumerate}[label={[\arabic*]},noitemsep,leftmargin=*]

\item Mirceski V, Skrzypek S, Stojanov L (2018). Square-wave voltammetry:
  a review on the recent progress. \textit{ChemTexts} 4:17.
  DOI:~10.1007/s40828-018-0073-0. \textbf{mirceski2018.pdf} -- surface-confined
  reversible SWV peak shape (sech$^2$ form), FWHM, quasi-reversible broadening,
  and peak area behaviour confirmed.

\item Savitzky A, Golay MJE (1964). Smoothing and differentiation of data
  by simplified least squares procedures. \textit{Anal Chem} 36:1627--1639.
  DOI:~10.1021/ac60214a047. \textbf{savitzky1964.pdf} -- convolution weights
  $\{c_k\}$ for polynomial smoothing (Table~1), noise reduction factor, and
  peak-preservation property confirmed.

\item Baker BR, Lai RY, Wood MS, Doctor EH, Heeger AJ, Plaxco KW (2006).
  An electronic, aptamer-based small-molecule sensor for the rapid,
  label-free detection of cocaine in adulterated samples and biological
  fluids. \textit{J Am Chem Soc} 128:3138--3139. DOI:~10.1021/ja056957p.
  \textbf{baker2006.pdf} -- SWV scan range and 2~mV step size confirmed;
  on-chip RE peak position reference confirmed.

\item Lubin AA, Lai RY, Baker BR, Heeger AJ, Plaxco KW (2010).
  Electrochemical, structure-switching biosensors: quantitative analysis
  of the usable dynamic range. \textit{Anal Chem} 78:5671--5677.
  DOI:~10.1021/ac0601417.
  \textbf{lubin2010.pdf} -- MB peak position, probe density range, and
  per-scan signal stability confirmed.

\item White RJ, Plaxco KW (2010). Exploiting binding-induced changes in
  probe flexibility for the optimization of electrochemical biosensors.
  \textit{Anal Chem} 82:73--76. DOI:~10.1021/ac9028595.
  \textbf{white2010.pdf} -- SWV frequency dependence of signal area;
  peak area robustness to frequency change compared to peak height
  confirmed.

\end{enumerate}

\end{document}
